\section{Inversion and threshold
certification}\label{tcl:inversion-and-threshold-certification}

\subsection{Inversion on the sequence
range}\label{tcl:inversion-on-the-sequence-range}

Let \(Y\) denote the target value. Define

\[
 \Phi(x)=2x(r-1+1/r)-x\log2,
 \qquad r=W(2x),
\]

\[
 h_0(r)=r^2/4+r/2+1+\tfrac12\log(1+r).
\]

Direct differentiation using \(r'=r/[x(1+r)]\) gives

\[
 \Phi'(x)=2r-\log2=:q(r),\qquad
 \Phi''(x)=\frac{2r}{x(1+r)}.
 \tag{C17}\label{tcl:eq:C17}
\]

From (C2),

\[
 \log A_P(n)=\Phi(n)-h_0(r)+a_{P,1}(r)/n+O_{P}(r^8/n^2).
 \tag{C18}\label{tcl:eq:C18}
\]

For large \(Y\), let \(x_0\) be the unique large real solution of
\(\Phi(x_0)=\log Y\) and let \(r_0=W(2x_0)\). If \(Y=A_P(N)\) for an
integer \(N\) tending to infinity, then

\[
 N=x_0+\frac{h_0(r_0)}{2r_0-\log2}
 +O_P(r_0^3/x_0).
 \tag{C19}\label{tcl:eq:C19}
\]

This is an absolute index error, not a relative one. To see the scale,
the leading displacement is \(O(r_0)\); inserting that displacement in
(C18) leaves \(O(r_0^4/x_0)\), which is divided by \(\Phi'\) asymptotic
to \(2r_0\).

One additional explicit correction is available. Write

\[
 d_0(r)=h_0(r)/q(r),\quad
 d_1(r)=\frac{\frac r{1+r}h_0'(r)d_0(r)
 -\frac r{1+r}d_0(r)^2-a_{P,1}(r)}{q(r)}.
\]

A Taylor expansion of (C18), now retaining its \(O(r^8/n^2)\) error,
gives

\[
 N=x_0+d_0(r_0)+d_1(r_0)/x_0+O_P(r_0^7/x_0^2).
 \tag{C20}\label{tcl:eq:C20}
\]

For details of the second correction, let $x=x_0$, $r=W(2x)$ and
$G(t)=\Phi(t)-h_0(W(2t))+a_{P,1}(W(2t))/t$.
For a displacement $d=O(r)$, finite Taylor expansion gives
\[
G(x+d)-\log Y=q(r)d-h_0(r)
 +\frac{\frac r{1+r}d^2-\frac r{1+r}h_0'(r)d+a_{P,1}(r)}{x}
 +O_P(r^8/x^2).
\]
Indeed $a_{P,1}=O(r^4)$ and its rational derivative is $O(r^3)$;
all higher Taylor terms are smaller than the displayed error.
Taking $d=d_0+d_1/x$ cancels the constant and $x^{-1}$ terms.
Equation (C18) at the true integer $N$, followed by the mean value theorem
for the explicit $G$ with $G'(t)\sim2W(2t)$, proves (C20).
The preliminary localization $N=x_0+O(r_0)$ follows first from (C18)
and $\Phi'(x)\asymp r_0$ in a fixed relative neighborhood of $x_0$.
No smooth continuation of the unknown error is used.

For arbitrary fixed \(J\), set

\[
 F_J(x)=\Phi(x)-h_0(W(2x))
 +\log\left(\sum_{j=0}^J a_{P,j}(W(2x))/x^j\right).
 \tag{C21}\label{tcl:eq:C21}
\]

The logarithm has positive argument eventually, and direct
differentiation of the finite rational formula shows
\(F_J'(x)\ge W(2x)\) eventually. If \(x_J\) is its large root
\(F_J(x_J)=\log Y\) and \(Y=A_P(N)\), then

\[
 |x_J-N|=O_{P,J}\!\left((1+r_N)^{4J+3}/N^{J+1}\right).
 \tag{C22}\label{tcl:eq:C22}
\]

No differentiation of the unknown asymptotic remainder is used: apply
(C2) only at the integer \(N\), and apply the mean value theorem to the
explicit smooth \(F_J\).

\subsection{Lambert seed and its
limitation}\label{tcl:lambert-seed-and-its-limitation}

Set

\[
 a=1+\tfrac12\log2,\qquad b=1-a^2/4,
\]

\[
 w=2W\!\left(\tfrac12\sqrt{\log Y}\,e^{-a/4}\right).
\]

The equation \(\Phi(x_0)=\log Y\) is exactly

\[
 e^{r_0}(r_0^2-ar_0+1)=\log Y.
\]

Writing \(r_0=w+a/2+\delta\) reduces it to

\[
 e^\delta((w+\delta)^2+b)=w^2,
\]

so

\[
 \delta=-\frac{b}{w(w+2)}+O(w^{-4}).
 \tag{C23}\label{tcl:eq:C23}
\]

This provides a useful initial seed for solving the exact scalar
equation for \(r_0\), then \(x_0=r_0e^{r_0}/2\). The defining Lambert
relation and branch convention are standard; see
\href{https://dlmf.nist.gov/4.13}{DLMF §4.13}.

A fixed finite truncation of this \(1/w\) seed does \textbf{not} retain
the fine absolute accuracy in (C19), (C20), or (C22). An \(r_0\) error
\(O(w^{-K})\) produces an \(x_0\) error \(O(x_0w^{-K})\), enormously
larger than the desired \(O(r_0^3/x_0)\). Refine the exact scalar
equation numerically to the needed precision, or keep the seed error
explicitly in the final error bound.

\subsection{Arbitrary thresholds require a true-sequence
enclosure}\label{tcl:arbitrary-thresholds-require-a-true-sequence-enclosure}

The asymptotic \(A_P(n)/H_n\to1\) and \(H_{n+1}/H_n\to\infty\) imply
\(A_P(n)\) is positive and strictly increasing eventually. For large
targets exceeding the maximum of the finite exceptional prefix, let
\(N_Y\) be the least integer with \(A_P(N_Y)\ge Y\).
If $n_*$ is an index beyond which positivity and monotonicity hold, this
least index lies in the increasing tail once
$Y>\max_{0\le k<n_*}A_P(k)$.
Equation (C22), as stated, concerns targets on the sequence range. For
arbitrary \(Y\), simply rounding a root of \(F_J\) does not
automatically certify \(N_Y\).

After the exceptional prefix has been excluded, a rigorous threshold
bracket requires either adjacent exact counts in the increasing tail, or a
proved bound

\[
 |\log A_P(k)-F_J(k)|\le E_J(k)
\]

with an explicit applicable constant and starting index. Integers
\(k_-\) and \(k_+\), both in that certified increasing tail, then certify

\[
 F_J(k_-)+E_J(k_-)<\log Y,\qquad
 F_J(k_+)-E_J(k_+)\ge\log Y
\]

and hence \(k_-<N_Y\le k_+\). A numerical bracket for the model root
alone does not provide those inequalities for the true sequence. The
theorem's existence of an \(O\)-constant is not an implemented
certified error bound.

For $V,U,L$ within the public exact range, a finite scan checks all earlier
indices instead of relying on an unimplemented monotonicity onset. Outside
that range, the present asymptotic theorem alone supplies no effective
certificate for a particular arbitrary threshold.

\begin{remark}[{[write]} The inverse and the transseries volume, 7~October 2026]\label{tcl:rem:transseries}
Compared statement by statement with the repository's transseries
volume~\cite{TSvol}; the volume's symbols are not used here.
\begin{enumerate}\raggedright
\item \emph{The growth.} $\log A_P(n)=2n\log n-2n\log\log n+O(n)$ (S18),
with a term $n\log n$, so the sequences are not of the exponential--power form
of \file{p0:def:model}; no instance of the volume's model is claimed.
\item \emph{The model root and the Lambert seed.} In the variable
$r_0=W(2x_0)$ the model equation $\Phi(x_0)=\log Y$ is exactly
$e^{r_0}(r_0^2-ar_0+1)=\log Y$ (checked symbolically), that is
$r_0+2\log r_0+\log(1-a/r_0+1/r_0^2)=\log\log Y$. This is literally the
volume's monomial--logarithmic equation
(\file{plt:def:lw-monomial-log-datum}) in the unknown $r_0$, with unit slope,
logarithmic coefficient~$2$, zero constant and a perturbation
$\log(1-at+t^2)$, $t=1/r_0$, of positive order with constant coefficients;
the three hypotheses of \file{plt:thm:lw-template} hold, so the formal
expansion of $r_0$ in the reciprocal of $\log\log Y$ is an exact instance of
the template. Its dominant block is solved exactly by the seed:
$w=2W(\frac12\sqrt{\log Y}\,e^{-a/4})$ is the solution of
$w+2\log w=\log\log Y-a/2$, an exact instance of \file{p0:thm:lambert-core}
(unit slope, positive logarithmic coefficient~$2$, principal branch), and
with $r_0=w+a/2+\delta$ the expansion~\eqref{tcl:eq:C23} is the first term of
the remaining correction, a series in~$1/w$ with constant coefficients.
\item \emph{The index.} The refinements~\eqref{tcl:eq:C19}, \eqref{tcl:eq:C20}
and~\eqref{tcl:eq:C22} move $x_0$ by $O(r_0)$, a relative change of order
$r_0e^{-r_0}$: below every order of the template's chart. So they are not
instances of \file{plt:thm:lw-template}; they are proved directly, at the
integer~$N$, by the mean value theorem for the explicit models.
[Independent check, 7~October 2026: the relative change is of order
$e^{-r_0}$, not $r_0e^{-r_0}$: the displacement $d_0(r_0)=h_0(r_0)/(2r_0-\log2)
\sim r_0/8$ divided by $x_0=r_0e^{r_0}/2$ is $\sim e^{-r_0}/4$. The printed
order is an upper bound, and the conclusion is unchanged.]
\item \emph{The thresholds.} $N_Y$ is the volume's staircase
(\file{p0:def:three-inverses}(1)) of $A_P$ on its increasing tail. The models
$\Phi-h_0$ and $F_J$ are smooth and increasing but do not take the values
$\log A_P(k)$ at the integers, so they are not admissible interpolations:
reading~\eqref{tcl:eq:C22} as $N=\lfloor x_J+\frac12\rfloor$ for large~$N$ is an
analogue of \file{p0:thm:staircase}(3), and the bracket $k_-<N_Y\le k_+$ of
Section~\ref{tcl:arbitrary-thresholds-require-a-true-sequence-enclosure},
proved directly at the integers from a true-sequence enclosure, is an
analogue of \file{p0:thm:staircase}(2), not an instance; the source's
caution that rounding a model root certifies nothing is the volume's
separation caveat.
\end{enumerate}
\end{remark}
